\subsection[Kollisionsprüfung gegenüber der Basis]{Kollisionsprüfung gegenüber der Basis \myAuthor{Johannes Decker}}\label{control/base_collision}
Im Rahmen des zweiten Meilensteins "MS2 Basic functionality and Inverse kinematics" des Control-Teams wurde das Thema Kollision des Roboters diskutiert. Zuerst wurde es für sinnvoll erachtet eine Kollisionsprüfung gegenüber der Basisebene $(z = 0)$ durchzuführen, da der Roboter für die Manipulation der Puzzleteile sich stehts nah an dieser bewegt. Es sollte dabei verhindert werden das der Roboter die Basisebene durchdringt und im Raum für $z \leq 0$ agiert, da Nutzung dieses Bewegungsraums die Grenzen eines realen Roboters überschreitet und keine sinnvolle Abstraktion eines realen Roboters darstellt. Da dieser Fall als potenziell mögliche Lösung der inversen Kinematik nicht zur Anwendung kommen darf, sollte eine Logik zur Kollisionsprüfung gegenüber der Basis entwickelt werden. 
Der gegebene Roboter ist ähnlich zu einem menschlichen Arm aufgebaut. Das heißt alle Segmente des Roboters folgen sequenziell aufeinander, sodass jedes Segment des Roboters mit maximal zwei Gelenken verbunden ist. Aus dem Aufbau des Roboters lässt sich also schließen das bei einer Kollision gegenüber der Basis, bzw. ein Eindringen in den $z \leq 0$ Raum, sich mindestens ein Gelenk des Roboters im $z \leq 0$ Raum befindet. Es muss also für jeden Bewegungsschritt des Roboters geprüft werden ob sich alle Gelenke im $z \geq 0$ Raum befinden. Dafür wird die Vorwärtskinematik für den Endpunkt des Bewegungsschritts berechnet. Anhand dieser kann für alle Gelenke des Roboters die Position in z-Richtung bestimmt und geprüft werden. Wenn der Wert der z-Koordinate eines jeden Gelenks des Roboters $ \geq 0$ ist, liegt keine Kollision gegenüber der Basis vor. Eine erste Version der Kollisionsprüfung gegenüber der Basis wurde durch die Funktion \texttt{base\_collision} in \texttt{src/control/pkg\_control/pkg\_control/robot.py} implementiert.

Im weiteren Verlauf des Meilensteins "MS2 Basic functionality and Inverse kinematics" viel jedoch auf das eine Kollisionsprüfung gegenüber der Basis aufgrund des Aufbaus des Roboters komplett ausgeschlossen werden kann. Aus Figure \ref{fig:rob_aufbau} kann abgelesen werden das die Länge des 2. Segmentes ("Oberarm") 25cm kürzer ist als die Länge des 1. Segments. Dadurch ist eine Kollision gegenüber der Basis für das Gelenk 3 und alle vorherigen Segmente und Gelenke des Roboters unmöglich. Ebenfalls können durch die Inverse Kinematik auch keine Konfigurationen der Gelenkwinkel mehr gefunden werden, in welchen der $z \leq 0$ Raum durchfahren wird. Aus diesem Grund wurde die Entwicklung und Implementierung der Funktion \texttt{base\_collision} nicht weiterverfolgt. Die Funktion findet auch in der finalen Implementierung keine Verwendung.

\begin{figure}[H]
    \centering
    \includegraphics[width=1.0\textwidth]{Kinematik_Arme.png}
    \caption{Aufbau des Roboters KUKA-KR6}
    \label{fig:rob_aufbau}
\end{figure}

